\section{Introduction}
\label{sec:introduction}

Shafarevich asked whether the universal cover of a complete algebraic
variety is holomorphically convex \citep{Shafarevich1974}.
We follow the formulation reproduced by
\citet[pp.~135--137]{LasellRamachandran1996}, who locate the question
on p.~407 of Shafarevich's book.
We use its usual smooth-projective formulation, called the
\emph{Shafarevich conjecture on holomorphic convexity}: the universal
cover of every smooth connected projective complex variety is
holomorphically convex. For a complex manifold $V$ and a nonempty compact set $K\subset V$, its \emph{holomorphic hull} is
\[
 \widehat K_V=\{v\in V: |g(v)|\leq\sup_K|g|
                    \text{ for every }g\in\mathcal O(V)\}.
\]
Holomorphic convexity means that every such hull is compact. We disprove the conjecture in complex dimension two.

\begin{theorem}\label{thm:main}
There is a smooth connected projective complex surface $X$ containing a connected nodal curve $Z_0$ with smooth rational irreducible components such that
\[
 \im\bigl(\pi_1(Z_0)\longrightarrow\pi_1(X)\bigr)
 \quad\text{is infinite}.
\]
In the simply connected universal cover of $X$, a connected component above $Z_0$ is a closed, connected, noncompact, locally finite union of compact rational curves. The holomorphic hull of any one of its points is noncompact. In particular, this universal cover is not holomorphically convex.
\end{theorem}

\subsection{Earlier approaches and the obstruction}

The conjecture connects projective geometry with the analytic geometry
of covering spaces. Gurjar and Shastri proved it for smooth projective
elliptic surfaces \citep[Theorem~2(ii)]{GurjarShastri1985}, and
Katzarkov proved it for smooth projective varieties with virtually
nilpotent fundamental group \citep{Katzarkov1997}.
Eyssidieux proved holomorphic convexity for the cover defined by the
intersection of the kernels of all reductive complex representations
of a fixed rank \citep{Eyssidieux2004}. Eyssidieux, Katzarkov, Pantev
and Ramachandran then established the universal-cover conclusion for
smooth projective varieties whose fundamental groups admit faithful
finite-dimensional complex representations
\citep{EyssidieuxKatzarkovPantevRamachandran2012}.
Campana, Claudon and Eyssidieux extended this linear Shafarevich
theorem to compact K\"ahler manifolds
\citep{CampanaClaudonEyssidieux2015}.

Recent results extend this theory beyond smooth projective varieties.
Deng and Yamanoi, with an appendix joint with Katzarkov, proved the
fixed-rank reductive-cover conclusion for normal projective varieties
\citep[Theorem~C(i)]{DengYamanoi2024}. Bakker, Brunebarbe and Tsimerman
proved that, for a connected normal algebraic space whose fundamental
group has a finite-dimensional complex representation with finite
kernel, the universal cover is a dense analytic-Zariski open subset
of a holomorphically convex complex space
\citep[Theorem~1.1]{BakkerBrunebarbeTsimerman2026}. This last conclusion
concerns a partial compactification of the universal cover.

The geometric reduction approach seeks to contract subvarieties for which
the fundamental group of the normalization has finite image in the ambient
fundamental group. Koll\'ar's rational Shafarevich maps
and Campana's meromorphic $\Gamma$-reduction address this problem through
sufficiently general points \citep[\S1]{Kollar1993}
\citep[Theorem~3.5 and Definition~3.8]{Campana1994}.
Such generic reductions do not by themselves provide a proper holomorphic
map from the entire universal cover to a Stein space.

The obstruction used in the present construction is classical.
Lasell and Ramachandran \citep[pp.~136--137]{LasellRamachandran1996}
and Katzarkov and Ramachandran
\citep[pp.~527--528]{KatzarkovRamachandran1998} explain why a
connected curve can obstruct holomorphic convexity when its normalized
components have finite fundamental-group image but the curve itself has
infinite image. A connected lift then has compact irreducible
components but is itself noncompact. Such a curve is called an
\emph{infinite Nori string}
\citep[arXiv version, Definition~3]{ColtoiuJoita2011}.
Every holomorphic function is constant on each compact component,
and connectedness makes these constants agree. A closed infinite Nori
string therefore cannot lie in a holomorphically convex space.
For our rational components the individual fundamental groups are
trivial; the decisive issue is the image of the loops in the incidence
graph of their whole nodal union.

Bogomolov and Katzarkov proposed constructions from degenerating
curves and quotients of their fundamental groups in which this
obstruction would occur if the resulting geometric quotient remained
infinite after all required relations
\citep[author manuscript, \S4.1, Lemma~4.2 and Conjecture~4.1]{BogomolovKatzarkov1998}.
That infinitude assertion remained conditional. Eyssidieux and Funar
subsequently proved holomorphic convexity for the associated
uniform-ramification construction from semistable families of curves
of genus at least two, outside explicit finite exceptional ranges
\citep[Theorems~1.2 and~6.18]{EyssidieuxFunar}.

The construction here follows the same broad passage from a family
of curves to a group quotient and then to a projective surface.
Its input is a marked genus-zero family, with both horizontal marked
sections and vertical boundary components. A \emph{meridian} is the
boundary of a small disk transverse to a divisor. We kill the meridians of
one distinguished fiber and impose high powers of the remaining
boundary meridians only after compactification. The argument must
establish both infinitude of this particular quotient and its
survival in the final surface. The weighted presentation proves the
first assertion; finite detection of entire local boundary-group
images, followed by transfer to the compact rational fiber, proves
the second. These are the steps developed below.

\subsection{Consequences and comparisons}

By the linear Shafarevich theorem, the fundamental group of the surface
in \Cref{thm:main} has no faithful finite-dimensional complex
representation. Two further comparisons concern restrictions on the
surface and on its universal cover.

A smooth compact K\"ahler manifold $M$ is \emph{special} in Campana's
sense if no holomorphic line bundle $L$ admits a nonzero map
$L\to\Omega_M^p$ with $\kappa(M,L)=p$ for any $1\le p\le\dim M$;
here $\kappa$ denotes Iitaka dimension
\citep[Theorem~2.27]{Campana2004}. The abelianity theorem for
special compact K\"ahler manifolds gives a virtually abelian
fundamental group, and hence a faithful finite-dimensional complex
representation. The compact K\"ahler linear theorem therefore makes
their ordinary universal covers holomorphically convex
\citep[Corollary~1.4]{OpenAIAbelianity2026}.
Consequently the surface in \Cref{thm:main} is not special.

A subset of a complex projective variety is \emph{semialgebraic}
if it is given in affine charts by finite Boolean combinations of
real polynomial equalities and inequalities. The ordinary universal
cover of our surface is not biholomorphic to a semialgebraic open
subset of a projective variety, with openness in the complex
topology: such a presentation would imply holomorphic convexity by
\citep[Corollary~8.1]{OpenAISemialgebraicCovers2026}, contrary to
\Cref{thm:main}.

\subsection{The construction in outline}

We begin with a family of spheres with $48$ moving punctures. The family comes from multiplication by $4$ on the elliptic curve with an automorphism of order three. Near one parameter value, the punctures collide in $24$ pairs. A double base change and one blowup at each collision produce an unmarked main sphere with $24$ attached spheres, each carrying two marked points. Finite covers of these components, branched only at those points, are again rational.

The main task is to retain an infinite fundamental-group image after filling in this fiber and all other missing fibers. Three features make this possible.

First, we study the moving points on the elliptic curve before taking its sign quotient. Three explicit double covers show that monodromy changes each generator only in higher degree. Deep elements of the free parameter group therefore impose relations of arbitrarily high degree. A surjective map of the actual based fiber groups transfers these equalities to the sphere family.

Second, the multiplication cover has deck group $(\Z/4)^2$. Successive differences under its two translations give a free basis with useful degrees. A calculation with the linear terms of the sign involution selects $24$ relations that impose all sign identifications on this $48$-generator group. These relations leave a strict margin in a weighted Golod--Shafarevich inequality. That margin accommodates the deep monodromy relations and large powers of the remaining boundary loops. The weighted infinitude method builds on \citet{GolodShafarevich1964}, \citet{Golod1964}, \citet{Vinberg1965}, \citet{Ershov2011}, and \citet{ErshovJaikinZapirain2013}. The group construction takes place in inverse limits of finite 2-groups, described in \Cref{sec:weighted}.

Third, we choose the finite parameter cover and compactify the marked family. Only then do we impose high powers of the remaining boundary meridians so that their images are finite. We then choose a finite quotient that is injective on each entire local boundary-group image. In the associated cover, this lets the infinite quotient survive normalization and resolution, including every exceptional divisor. The resulting smooth projective surface retains a rational special fiber. A neighborhood deformation retraction and properness transfer the infinite image of a nearby fiber to a connected component of that rational fiber.

The two reusable mechanisms are the economical involution presentation in \Cref{prop:involution} and the finite-local-group argument in \Cref{lem:local-detection}. The former retains enough generators relative to relations for the weighted infinitude criterion; the latter extends a homomorphism from the fundamental group of a divisor complement across a smooth projective compactification of a finite cover.

We give the weighted algebra in \Cref{sec:weighted}, construct and compare the actual fiber groups in \Cref{sec:family}, and analyze the involution in \Cref{sec:involution}. \Cref{sec:quotient} constructs the infinite quotient and the compactified marked family. \Cref{sec:surface} completes the surface and its rational curve.

\subsection{Why the rational curve is enough}

The following elementary form of the compact-curve obstruction makes the last step explicit. It is the rational-curve case of the mechanism described in \citet[pp.~136--137]{LasellRamachandran1996}, \citet[pp.~527--528]{KatzarkovRamachandran1998}, and \citet[author manuscript, \S4.1]{BogomolovKatzarkov1998}.

\begin{lemma}\label{lem:rational-hull}
Let $X$ be a connected complex manifold and let $Z\subset X$ be a compact connected nodal curve whose irreducible components are smooth rational curves. If the image of $\pi_1(Z)\to\pi_1(X)$ is infinite, then every connected component $W$ of the inverse image of $Z$ in the simply connected universal cover $\widetilde X$ is closed and noncompact. It is a locally finite union of compact rational curves, and every holomorphic function on $\widetilde X$ is constant on $W$.
Consequently $\widehat{\{w\}}_{\widetilde X}$ is noncompact for every $w\in W$.
\end{lemma}

\begin{proof}
The restricted map $W\to Z$ is the connected covering corresponding to
\[
 \ker\bigl(\pi_1(Z)\longrightarrow\pi_1(X)\bigr).
\]
Thus it has infinitely many sheets. A fiber is a closed discrete infinite subset, so $W$ cannot be compact. It is closed in $\widetilde X$, since it is a connected component of the closed inverse image of $Z$.

Each rational component of $Z$ is simply connected. Its inverse image therefore consists of compact rational curves mapping isomorphically onto it. Choose an evenly covered neighborhood in $X$ of each point of $Z$, small enough that it meets $Z$ in one smooth branch or in the two branches of a node. Each lifted neighborhood meets at most one or two lifted components. These neighborhoods show that the collection is locally finite in the ambient manifold $\widetilde X$; points outside the closed inverse image have a neighborhood meeting none of it. The connected space $W$ is locally path connected; a path in it has compact image and meets only finitely many of these curves. The maximum principle makes a holomorphic function constant on each compact curve, and the constants agree at their intersections. They therefore agree throughout $W$.

For $w\in W$, this proves $W\subset\widehat{\{w\}}_{\widetilde X}$. A compact hull cannot contain the closed noncompact set $W$.
\end{proof}
